\section{Conclusion}
\label{sec:conclusion}

We investigated how \mistral internally represents French numbers, with a focus on the vigesimal counting system used for numbers 70--99. Using circular digit probes across four input formats, we found that French number word representations encode base-10 digits at 98.0\% accuracy---comparable to English and significantly higher than digit strings. The vigesimal structure leaves localized fingerprints: errors on the \emph{soixante-dix} (70s) range reveal partial encoding of the arithmetic decomposition (``soixante'' $\to$ 60 rather than the composed value 70+). Most strikingly, number word representations peak at early layers (5--6), while digit string representations peak at the final layer (31), suggesting distinct processing pathways for semantic and tokenized numerical inputs.

Our work demonstrates that distributional learning is sufficient to capture the implicit arithmetic of complex counting systems. The French vigesimal system---where 98 is encoded as $4\times20+10+8$ in the surface form---does not prevent the model from building clean base-10 digit representations.

Several directions remain open. Testing additional models (multilingual LLMs, models at different scales) would establish generality. Causal interventions could confirm that these representations are functionally used. Token-by-token analysis of compound vigesimal words could trace how arithmetic composition unfolds across the input. Finally, extending this analysis to other complex counting systems---Danish (\emph{halvfjerds} = ``half-fourth-times-twenty'' for 70), traditional Japanese, or Basque---would further illuminate the relationship between linguistic structure and internal numeric representations.
